
\section{Introduction}
\label{sec:intro}

Long-horizon agents increasingly act over state assembled from memory, retrieval, tools,
logs, and prior decisions. The downstream decision policy rarely sees the entire underlying
world state. Instead, it sees a \emph{decision-visible evidence surface}: a finite set of
records selected by upstream retrieval and filtering stages.

This creates a safety problem that ordinary relevance metrics do not fully capture. A
policy can be conservative when missing information is itself visible. If an approval
requirement remains visible but the approval record is absent, the policy can request
approval. If a stale record is visible without a current replacement, the policy can
defer. If two claims conflict and no resolution record is present, the policy can clarify.
We call these \emph{detectable gaps}.

A different failure occurs when evidence loss removes both a required item and the only
cue that anything is missing. Consider an earlier instruction to execute an action followed
by a later override cancelling it. If retrieval drops the cancellation, the remaining
surface can contain a coherent older instruction plus readiness evidence. A downstream
policy that checks only what it can see has no local signal that a later override exists.
We call this a \emph{silent evidence omission}.

\begin{figure}[t]
\centering
\resizebox{\linewidth}{!}{% Conceptual figure: no empirical data.
\begin{tikzpicture}[
  node distance=6mm and 7mm,
  every node/.style={font=\small},
  box/.style={draw, rounded corners=2pt, align=center, inner sep=5pt, minimum width=31mm},
  decision/.style={draw, rounded corners=2pt, align=center, inner sep=5pt, minimum width=32mm, fill=black!4},
  arrow/.style={-Latex, thick}
]
\node[box] (world) {Underlying world\\and source truth};
\node[box, below=of world] (retrieval) {Evidence acquisition\\and retrieval};
\node[box, below=of retrieval] (surface) {Decision-visible\\evidence surface};

\node[decision, below left=10mm and 7mm of surface] (visible) {Obligation uncovered\\and gap cue visible};
\node[decision, below right=10mm and 7mm of surface] (silent) {Obligation uncovered\\but no gap cue visible};

\node[box, below=of visible] (fallback) {CLARIFY / DEFER /\\REQUEST APPROVAL};
\node[box, below=of silent] (unsafe) {Older / local state\\still looks actionable};

\node[box, below=of unsafe] (risk) {Unsupported-decision\\risk};
\node[decision, right=14mm of risk] (coverage) {Independent\\coverage-complete signal};
\node[box, below=of coverage] (safe2) {Coverage-aware\\fallback};

\draw[arrow] (world) -- (retrieval);
\draw[arrow] (retrieval) -- (surface);
\draw[arrow] (surface) -- (visible);
\draw[arrow] (surface) -- (silent);
\draw[arrow] (visible) -- (fallback);
\draw[arrow] (silent) -- (unsafe);
\draw[arrow] (unsafe) -- (risk);
\draw[arrow] (risk) -- node[above, font=\scriptsize]{if incompleteness is exposed} (coverage);
\draw[arrow] (coverage) -- (safe2);
\end{tikzpicture}
}
\caption{\textbf{Missing evidence is not a single failure mode.} A downstream policy can
react safely when an uncovered obligation leaves a visible gap cue. A silent omission
removes the required evidence without leaving such a cue. The coverage-aware condition
tests a separate completeness signal without revealing the missing content.}
\label{fig:mechanism}
\end{figure}

This failure is distinct from two nearby problems. STALE studies whether agents revise
state correctly when later evidence is available but implicitly invalidates older memory
\citep{chao2026stale}. Revocation-enforcement work studies the opposite read-path error:
an already revoked record remains retrievable and authoritative
\citep{shen2026revoked}. Our intervention removes the \emph{superseding valid record}
from the decision-visible surface, so the older record can remain locally coherent without
any observable revocation cue.

We ask three questions:
\begin{enumerate}
  \item \textbf{Causal visibility:} does hiding formally required evidence, while holding
  task and policy fixed, increase unsupported decisive behavior?
  \item \textbf{Failure mechanism:} are all missing obligations equally dangerous, or are
  failures concentrated where absence is silent relative to visible-gap checks?
  \item \textbf{Coverage signaling:} can a signal that required obligations are incomplete
  prevent unsupported decisions without revealing the missing content itself?
\end{enumerate}

We study these questions in a preregistered, formally specified synthetic benchmark with
180 scenarios spanning six task families and 12 domains. Each scenario contains a
machine-readable source truth, explicit evidence obligations, rendered evidence, a frozen
decision-safety contract, and paired FULL/HIDDEN treatments. Two separately implemented
formal validators must agree before a sample is admitted.

The confirmatory experiment uses one frozen deterministic downstream policy and six
conditions. C0 exposes all formal evidence. C1 removes exactly one designated mandatory
item. C2--C4 use frozen lexical, hashed-dense, and hybrid top-$k$ retrieval. C5 reuses the
\emph{identical C4 retrieval surface} and adds only a boolean
\texttt{coverage\_complete} signal.

The principal result is not that all missing evidence is dangerous. The paired C0$\to$C1
intervention raises UDR from 0\% to 16.67\%, but all 30 direct-treatment failures occur in
the override/invalidation family. In the other five families, the omission leaves a
policy-visible structural cue and the frozen policy falls back conservatively. Under
ordinary top-$k$ retrieval, however, unsupported decisions also appear in constraint,
authority, and conflict families because retrieval can remove combinations of evidence
or the cue-bearing record itself. Silentness is therefore a property of the
\emph{resulting surface relative to the policy}, not merely a task-family label.

Our contributions are:
\begin{itemize}
  \item a policy-relative formalization of \emph{silent evidence omission}, distinct from
  stale-state resolution and from retrieval that returns explicitly revoked records;
  \item a preregistered paired intervention that holds world state and downstream policy
  fixed while changing mandatory-evidence visibility;
  \item a boundary characterization showing that direct hiding is safe when a visible gap
  remains, but unsafe for superseding override/invalidation evidence whose omission leaves
  an apparently actionable surface;
  \item evidence that aggregate mandatory-fact recall does not induce a strict monotonic
  safety ordering across the frozen retrievers; and
  \item a same-surface mechanism test showing that obligation-completeness information is
  sufficient, under the formal benchmark, to remove the measured unsupported decisions
  without suppressing decisions on complete-coverage cases.
\end{itemize}

Our claims are deliberately narrow. The benchmark is synthetic and formal. It does not
establish human-preference validity, real-world failure prevalence, deployment safety, or
a universal theorem about agents.
